\documentclass[10pt, a4paper]{article}

% Packages
\usepackage[utf8]{inputenc}
\usepackage[T1]{fontenc}
\usepackage[margin=1in]{geometry}
\usepackage{xcolor}
\usepackage{hyperref}
\usepackage{fontawesome5}
\usepackage{charter}
\usepackage{parskip}

% Define Primary Color - Siemens-compatible Navy
\definecolor{primarycolor}{RGB}{0, 51, 102}

% Hyperref Setup
\hypersetup{
    colorlinks=true,
    linkcolor=primarycolor,
    urlcolor=primarycolor,
}

\begin{document}

% --- SENDER INFO ---
\begin{flushright}
    \textbf{\large Kartavya Niraj Dikshit} \\
    \small
    Halbmondstraße 7, 91054 Erlangen \\
    (+49) 15510940829 \\
    \href{mailto:kartavya.dikshit@gmail.com}{kartavya.dikshit@gmail.com} \\
    \href{http://www.linkedin.com/in/kartavya-dikshit}{linkedin.com/in/kartavya-dikshit} \\
    \href{https://github.com/KartavyaDikshit}{github.com/KartavyaDikshit}
\end{flushright}

\vspace{20pt}

% --- RECIPIENT INFO ---
\begin{flushleft}
    Siemens AG \\
    Siemens Advanta - Recruiting Team \\
    Otto-Hahn-Ring 6 \\
    81739 München, Deutschland
\end{flushleft}

\vspace{20pt}

% --- DATE & SUBJECT ---
\begin{flushright}
    Erlangen, 19. März 2026
\end{flushright}

\vspace{10pt}
\textbf{\large Bewerbung als Werkstudent (f/m/d) Advanced Analytics and Artificial Intelligence} \\
\textbf{Job-ID: 497378}

\vspace{20pt}

% --- SALUTATION ---
Sehr geehrte Damen und Herren,

Siemens Advanta ist weltweit führend darin, komplexe technologische Visionen in reale industrielle Wertschöpfung zu transformieren. Als Masterstudent der Data Science an der Friedrich-Alexander-Universität Erlangen-Nürnberg mit einer klaren Spezialisierung auf Artificial Intelligence und Machine Learning suche ich eine Herausforderung, bei der ich meine akademische Tiefe mit meinen praktischen Erfahrungen in Cloud-Architekturen kombinieren kann. Die ausgeschriebene Position in München bietet genau diese Schnittstelle.

In meinem aktuellen Studium an der FAU konzentriere ich mich auf Deep Learning, Scalable Data Engineering und Reinforcement Learning. Mein Ziel ist es nicht nur, Algorithmen zu entwickeln, sondern diese in produktionsreife, skalierbare Systeme zu überführen. Meine bisherigen Stationen als AI-Engineer haben mir gezeigt, dass der Erfolg von KI-Projekten maßgeblich von der Qualität der Datenpipelines und der nahtlosen Cloud-Integration abhängt. Mit fundierten Kenntnissen in Python, SQL und praktischer Erfahrung in Snowflake- sowie Azure-Umgebungen bringe ich die technischen Werkzeuge mit, um Ihre Teams bei der Extraktion und Transformation komplexer Datensätze sofort wirksam zu unterstützen.

Besonders reizt mich an dieser Rolle bei Siemens Advanta die Kombination aus technischer Implementierung und der Kommunikation von Erkenntnissen. In meinen bisherigen Projekten konnte ich bereits beweisen, dass ich komplexe analytische Modelle in interaktive Dashboards übersetzen kann, die für Stakeholder echte Entscheidungshilfen bieten. Zudem bringe ich Erfahrung in der Durchführung von Workshops und der Erstellung von technischen Schulungsunterlagen mit – eine Kompetenz, die ich gerne einsetzen möchte, um den Wissensaustausch innerhalb Ihres Teams zu fördern.

Ich bin ein analytischer Denker, der agiles Arbeiten und kollaborative Entwicklungsworkflows mit Git schätzt. Die Möglichkeit, 15 bis 20 Stunden pro Woche an zukunftsweisenden Projekten der Siemens Advanta mitzuwirken, motiviert mich sehr. 

Gerne überzeuge ich Sie in einem persönlichen Gespräch von meiner Eignung und meiner Begeisterung für die digitale Transformation bei Siemens.

Mit freundlichen Grüßen,

\vspace{20pt}
Kartavya Niraj Dikshit

\end{document}